\section{The project-local Math layer}
\label{sec:math-layer}

The library has a project-local \TT{Math/} namespace between
mathlib and the game-theoretic development.  It packages
mathematical infrastructure that is needed repeatedly by the game
library but is not itself game theory.  A declaration belongs here
when it can be stated without players, utilities, equilibria,
strategies, or mechanisms.  Once those words become part of the
statement, the result belongs in \TT{GameTheory/}.

This separation keeps the kernel-game API small.  The game-theory
files can speak in terms of expected utility, mixed profiles, traces,
or information states, while the reusable algebra about PMFs,
products, couplings, finite traces, DAGs, denominator clearing, and
online regret remains available to every language layer.

\subsection{Discrete probability infrastructure}
\label{sec:math-layer:probability}

The largest part of \TT{Math/} is the discrete-probability layer.
\TT{Math.Probability} introduces the library's notation for kernels
and expected values over mathlib's \TT{PMF}.  It contains the finite
sum wrappers used throughout the paper and the bounded-tsum lemmas
that let expected utility remain meaningful on infinite outcome
carriers when utilities are bounded.  \TT{Math.ProbabilityMassFunction}
adds pushforward, bind-congruence, support, conditioning, and
monotonicity lemmas for PMFs.

Product distributions are factored into \TT{Math.PMFProduct}.  The
basic operation is \TT{pmfPi}, the independent product of finitely
many player-indexed PMFs.  The companion files record the algebra
needed by mixed extensions and information-model proofs:\ coordinate
projection, coordinate update, bind/factorization laws,
conditioning, and independence from a coordinate.  This is why the
same \TT{pmfPi} facts can support normal-form mixed extensions,
correlation regimes, observation models, MAID compilation, and
FOSG/MultiRound execution.

The probabilistic relation layer is \TT{Math.Coupling}.  It defines
relation lifting for PMFs by a joint distribution with prescribed
marginals and support contained in the relation.  This is not a game
notion, but it is the right vocabulary for bisimulation-style
arguments:\ equality of laws, projection through a view, and
one-step relational simulation are all special cases.  The companion
\TT{Math.NondeterministicRefinement} file lifts the same discipline
from single laws to families of laws indexed by an unweighted
resolution type.

\subsection{Runs, traces, and finite process closure}
\label{sec:math-layer:runs}

Several language compilers need to iterate stochastic transitions.
\TT{Math.PMFIter} provides state-only Kleisli iteration:\ starting
from a state and a transition kernel, it returns the law of the state
after a fixed number of steps.  \TT{Math.TraceRun} records the entire
finite state trace instead of only the final state.  The
parameterized version, \TT{Math.ParameterizedChain}, proves the
tower-property calculation used by Kuhn-style arguments:\ averaging
over hidden pure strategies can be represented by a single sequential
chain with local conditioning.

\TT{Math.OutcomeClosure} packages a common finite-process argument.
If a continuation value is preserved by every nonterminal step, and
the process has a rank that decreases on supported transitions, then
running the stopped process for enough fuel and projecting terminal
states recovers the initial value.  FOSG outcome closure uses this
as a game-neutral lemma; the game-specific file only supplies the
state type, terminal predicate, transition kernel, rank, and
observation map.

\subsection{Finite combinatorics and ordered structure}
\label{sec:math-layer:combinatorics}

The language layer also uses finite combinatorics that is independent
of payoffs.  \TT{Math.DAG} formalizes acyclic binary relations and
topological orders, including the equivalence between acyclicity and
existence of a topological order for finite predecessor-set DAGs.
The MAID syntax uses this to separate graph well-formedness from
strategic interpretation.  \TT{Math.BoundedLists} enumerates lists of
bounded length over a finite alphabet, a small combinatorial tool
used by finite-horizon trace and history constructions.

\TT{Math.Reindex} contains generic finite-sum and \TT{ENNReal}
tsum reindexing lemmas.  \TT{Math.SimplexApproximation} contains
the denominator-clearing estimates for finite probability vectors.
Those estimates are the arithmetic substrate for turning a feasible
mixed payoff vector into a finite periodic path in the discounted
folk-theorem development.

\TT{Math.Fintype}, \TT{Math.Fin}, and
\TT{Math/Fintype/Transport.lean} package finite-cardinality and
\TT{Fin n} transport lemmas used to move theorem statements between
arbitrary finite carrier types and canonical \TT{Fin} indices.  These
are intentionally game-free: their statements mention only functions,
subtypes, finite cardinalities, and finite sums.

\subsection{Analysis, optimization, and learning}
\label{sec:math-layer:analysis}

Some infrastructure is analytic rather than combinatorial.
\TT{Math.OptimizationLocalGlobal} states a small abstract vocabulary
for fixed points, no-improvement maps, and local optimality.  The
mixed-Nash proof uses this style to isolate the algebraic step that
turns a fixed point of the Nash map into a no-profitable-deviation
profile.

\TT{Math.OnlineLearning} is the game-free online-learning component.
Its multiplicative-weights file defines the distribution updated by
exponential weights and proves the finite-horizon external-regret
bound.  The game-theoretic self-play theorem imports that result and
then identifies each player's online losses with external-regret
gains in the kernel game.  This prevents the no-regret proof from
being entangled with equilibrium terminology.

\TT{Math.SchauderFixedPoint} proves a Schauder fixed-point theorem
for nonempty compact convex subsets of real normed spaces, using the
standalone Brouwer theorem through a finite-net approximation.  The
main mixed-Nash theorem uses the finite-dimensional Brouwer wrapper,
not Schauder.  Schauder is included because it marks the mathematical
shape of a compact-strategy extension:\ the missing layer is the
game-theoretic compactness and continuity API, not the abstract
fixed-point theorem.

\TT{Math.FixedPoint.KKM} proves the finite
Knaster--Kuratowski--Mazurkiewicz covering lemma~\citep{KKM1929} for
closed covers of the standard simplex, plus an open-cover variant.  It
is used by the Stromquist fair-division proof, but the theorem itself
is stated without agents, utilities, or cake-cutting language:

\begin{theorem}[KKM covering lemma; \TT{kkm\_closed\_cover}]
\label{thm:kkm}
Let $F_i$ be closed subsets of the standard simplex indexed by
$i : \mathrm{Fin}\,n$.  If every face of the simplex is covered by the
sets whose indices belong to that face, then all the $F_i$ have a
common point.
\end{theorem}

\TT{Math.MeasureTheory.UnitInterval} supplies the measure-theoretic
unit-interval facts consumed by divisible fair division: pushforward of
a non-atomic measure along the subtype embedding, continuity of the
cumulative distribution function, interval-measure identities, and the
existence of cuts realizing a target measure value in $[0,1]$.

\subsection{Epistemic and semantic utilities}
\label{sec:math-layer:semantic}

Finally, \TT{Math.Knowledge} records the S5 knowledge operator
induced by a view function.  It is deliberately stated without
players, priors, or payoffs:\ indistinguishability is equality of
views, knowledge is truth on the equivalence class, and coarsening is
postcomposition by a projection.  The common-knowledge files in
\TT{Concepts/Knowledge/} add the game-theoretic and probabilistic
structure on top.

The same pattern appears throughout the Math layer.  It supplies
small, reusable mathematical contracts; the game library supplies
the interpretations.  That boundary is one reason the same semantic
core can support normal-form games, extensive-form games,
information models, stochastic games, mechanisms, auctions,
cooperative games, and learning dynamics without re-proving the
underlying probability and finite-process algebra in each setting.

\leancorr{\TT{Math/Probability.lean},
\TT{Math/ProbabilityMassFunction.lean}, \TT{Math/PMFProduct/},
\TT{Math/Coupling.lean}, \TT{Math/NondeterministicRefinement.lean},
\TT{Math/PMFIter.lean}, \TT{Math/TraceRun.lean},
\TT{Math/ParameterizedChain.lean}, \TT{Math/OutcomeClosure.lean},
\TT{Math/DAG.lean}, \TT{Math/BoundedLists.lean},
\TT{Math/Reindex.lean}, \TT{Math/SimplexApproximation.lean},
\TT{Math/Fintype.lean}, \TT{Math/Fintype/Transport.lean},
\TT{Math/Fin.lean}, \TT{Math/FixedPoint/KKM.lean},
\TT{Math/MeasureTheory/UnitInterval.lean},
\TT{Math/OptimizationLocalGlobal.lean}, \TT{Math/OnlineLearning/},
\TT{Math/SchauderFixedPoint.lean}, and \TT{Math/Knowledge.lean}.}
